\documentclass[11pt]{article}
\usepackage[margin=1in]{geometry}
\usepackage{amsmath, amssymb, amsfonts, bm}
\usepackage{array}
\usepackage{booktabs}
\usepackage{enumitem}
\usepackage{graphicx}
\usepackage{setspace}

\setlength{\parindent}{0pt}
\setlength{\parskip}{0.6em}

% --- Title Page ---
\begin{document}
\begin{titlepage}
    \begin{center}
        \vspace*{2cm}
        \Huge\textbf{Assignment 4} \\
        \vspace{0.25cm}
        \Large\textit{SYSC5108} \\
        \vspace{2cm}
        \Large Jesse Levine (101185127) \\
        \vfill
    \end{center}
\end{titlepage}

\newpage
\pagenumbering{arabic}

\section*{1}

Given:
\[
\bm{x}_t=
\begin{bmatrix}
1\\
0.5
\end{bmatrix},
\qquad
\bm{h}_{t-1}=
\begin{bmatrix}
0.05\\
0.2\\
0.15
\end{bmatrix},
\qquad
\bm{b}_h=
\begin{bmatrix}
0.1\\
0.1\\
0.1
\end{bmatrix},
\qquad
\bm{b}_y=
\begin{bmatrix}
0.1\\
0.1
\end{bmatrix}
\]

and the weight matrices:
\[
W_{xh}=
\begin{bmatrix}
-0.07 & 0.05\\
-0.04 & 0.10\\
-0.05 & -0.05
\end{bmatrix},
\quad
W_{hh}=
\begin{bmatrix}
-0.22 & -0.10 & 0.05\\
-0.04 & -0.09 & -0.06\\
-0.09 & 0 & -0.08
\end{bmatrix},
\quad
W_{hy}=
\begin{bmatrix}
0.06 & -0.01 & -0.18\\
-0.06 & 0.13 & 0.14
\end{bmatrix}.
\]

\begin{enumerate}[label=\alph*)]
\item Hidden-layer pre-activation:
\[
\bm{z}_h = W_{xh}\bm{x}_t + W_{hh}\bm{h}_{t-1} + \bm{b}_h.
\]

First,
\[
W_{xh}\bm{x}_t=
\begin{bmatrix}
-0.07(1)+0.05(0.5)\\
-0.04(1)+0.10(0.5)\\
-0.05(1)-0.05(0.5)
\end{bmatrix}
=
\begin{bmatrix}
-0.045\\
0.010\\
-0.075
\end{bmatrix},
\]
and
\[
W_{hh}\bm{h}_{t-1}=
\begin{bmatrix}
-0.22(0.05)-0.10(0.2)+0.05(0.15)\\
-0.04(0.05)-0.09(0.2)-0.06(0.15)\\
-0.09(0.05)+0(0.2)-0.08(0.15)
\end{bmatrix}
=
\begin{bmatrix}
-0.0235\\
-0.0290\\
-0.0165
\end{bmatrix}.
\]

Therefore,
\[
\bm{z}_h=
\begin{bmatrix}
0.0315\\
0.0810\\
0.0085
\end{bmatrix}.
\]

Since the hidden units use ReLU,
\[
\bm{h}_t=\operatorname{ReLU}(\bm{z}_h)=
\begin{bmatrix}
0.0315\\
0.0810\\
0.0085
\end{bmatrix}.
\]

Now compute the output pre-activation:
\[
\bm{z}_y=W_{hy}\bm{h}_t+\bm{b}_y
=
\begin{bmatrix}
0.06 & -0.01 & -0.18\\
-0.06 & 0.13 & 0.14
\end{bmatrix}
\begin{bmatrix}
0.0315\\
0.0810\\
0.0085
\end{bmatrix}
+
\begin{bmatrix}
0.1\\
0.1
\end{bmatrix}
=
\begin{bmatrix}
0.09955\\
0.10983
\end{bmatrix}.
\]

Again using ReLU at the output,
\[
\bm{y}_t=\operatorname{ReLU}(\bm{z}_y)=
\begin{bmatrix}
0.09955\\
0.10983
\end{bmatrix}.
\]

Therefore the answer is:
\[
\boxed{
\bm{y}_t=
\begin{bmatrix}
0.09955\\
0.10983
\end{bmatrix}}
\]

\item Let the target be
\[
\bm{t}=
\begin{bmatrix}
0\\
1
\end{bmatrix}.
\]

Using squared error and the same convention as earlier assignments,
\[
\delta_k=(y_k-t_k)\operatorname{ReLU}'(z_k).
\]

Since all hidden and output pre-activations are positive, every ReLU derivative is \(1\).
Thus the output-layer deltas are
\[
\delta_6 = 0.09955-0 = 0.09955,
\qquad
\delta_7 = 0.10983-1 = -0.89017.
\]

For the hidden layer,
\[
\bm{\delta}_h = (W_{hy}^{\top}\bm{\delta}_y)\odot \operatorname{ReLU}'(\bm{z}_h).
\]
So,
\[
W_{hy}^{\top}\bm{\delta}_y=
\begin{bmatrix}
0.06 & -0.06\\
-0.01 & 0.13\\
-0.18 & 0.14
\end{bmatrix}
\begin{bmatrix}
0.09955\\
-0.89017
\end{bmatrix}
=
\begin{bmatrix}
0.0593832\\
-0.1167176\\
-0.1425428
\end{bmatrix}.
\]

Therefore the hidden-layer deltas are
\[
\delta_3=0.0593832,\qquad
\delta_4=-0.1167176,\qquad
\delta_5=-0.1425428.
\]

Hence the neuron deltas are
\[
\boxed{
\delta_3=0.05938,\;
\delta_4=-0.11672,\;
\delta_5=-0.14254,\;
\delta_6=0.09955,\;
\delta_7=-0.89017.}
\]
\end{enumerate}

\section*{2}

From Slide 10 of Chapter 6,
\[
P=
\begin{bmatrix}
p & 1-p\\
0 & 1
\end{bmatrix},
\]
with the initial state at \(n=0\) equal to state \(0\), so
\[
\bm{\pi}(0)=
\begin{bmatrix}
1 & 0
\end{bmatrix}.
\]

\begin{enumerate}[label=\alph*)]
\item At time \(n\),
\[
\pi_0(n)=p^n,
\qquad
\pi_1(n)=1-p^n.
\]

\item To revisit state \(0\), the chain must stay in state \(0\) in the first step. After reaching
state \(1\), it never returns. Therefore,
\[
\boxed{\Pr(\text{revisit state }0)=p.}
\]

\item For \(0 \le p < 1\),
\[
\bar{\pi}_0=\lim_{N\to\infty}\frac{1}{N}\sum_{n=0}^{N-1}p^n=0,
\qquad
\bar{\pi}_1=\lim_{N\to\infty}\frac{1}{N}\sum_{n=0}^{N-1}(1-p^n)=1.
\]

Hence,
\[
\boxed{\bar{\pi}_0=0,\qquad \bar{\pi}_1=1.}
\]

\item Before the message is sent successfully, the sender can only be in state \(0\). Hence
\[
\boxed{\Pr(\text{sender is in state }0\mid \text{message not yet sent})=1.}
\]
\end{enumerate}

\section*{3}

The transition matrix is
\[
P=
\begin{bmatrix}
p & 1-p\\
1-p & p
\end{bmatrix}.
\]

\begin{enumerate}[label=\alph*)]
\item With
\[
Q_0=
\begin{bmatrix}
0.2\\
0.8
\end{bmatrix},
\]
the marginal distribution at time \(1\) is
\[
Q_1=PQ_0=
\begin{bmatrix}
p & 1-p\\
1-p & p
\end{bmatrix}
\begin{bmatrix}
0.2\\
0.8
\end{bmatrix}
=
\begin{bmatrix}
0.8-0.6p\\
0.2+0.6p
\end{bmatrix}.
\]

\item Define
\[
A=\frac{1}{2}
\begin{bmatrix}
1 & 1\\
1 & 1
\end{bmatrix},
\qquad
B=\frac{1}{2}
\begin{bmatrix}
1 & -1\\
-1 & 1
\end{bmatrix}.
\]
Then
\[
P=A+(2p-1)B.
\]

Also,
\[
A^2=A,\qquad B^2=B,\qquad AB=BA=0.
\]

Therefore,
\[
P^{(n)} = \left(A+(2p-1)B\right)^n = A+(2p-1)^nB.
\]
Substituting \(A\) and \(B\),
\[
P^{(n)}=
\begin{bmatrix}
\frac{1}{2}+\frac{(2p-1)^n}{2} & \frac{1}{2}-\frac{(2p-1)^n}{2}\\[4pt]
\frac{1}{2}-\frac{(2p-1)^n}{2} & \frac{1}{2}+\frac{(2p-1)^n}{2}
\end{bmatrix},
\]
which is exactly the required result.

\item If \(0<p<1\), then \(|2p-1|<1\), so
\[
\lim_{n\to\infty}(2p-1)^n=0.
\]
Hence the limiting probabilities are
\[
\boxed{
\lim_{n\to\infty}Q_n=
\begin{bmatrix}
\tfrac{1}{2}\\
\tfrac{1}{2}
\end{bmatrix}.}
\]

Special cases:
\[
p=1 \Rightarrow P=I \text{ and } Q_n=Q_0,
\qquad
p=0 \Rightarrow P^n \text{ oscillates and no limit exists.}
\]

\item For \(0<p<1\),
\[
\boxed{
\lim_{n\to\infty}P^{(n)}=
\begin{bmatrix}
\tfrac{1}{2} & \tfrac{1}{2}\\
\tfrac{1}{2} & \tfrac{1}{2}
\end{bmatrix}.}
\]

For \(p=1\), the limit is \(I\). For \(p=0\), the limit does not exist.
\end{enumerate}

\section*{4}

Let the original reward at time \(t+1\) be \(R_{t+1}\), and define a new reward
\[
R'_{t+1}=R_{t+1}+c.
\]

Under any policy \(\pi\),
\[
v'_{\pi}(s)
=
\mathbb{E}_{\pi}\!\left[\sum_{k=0}^{\infty}\gamma^k R'_{t+k+1}\,\middle|\,S_t=s\right]
\]
\[
=
\mathbb{E}_{\pi}\!\left[\sum_{k=0}^{\infty}\gamma^k (R_{t+k+1}+c)\,\middle|\,S_t=s\right]
\]
\[
=
\mathbb{E}_{\pi}\!\left[\sum_{k=0}^{\infty}\gamma^k R_{t+k+1}\,\middle|\,S_t=s\right]
+
c\sum_{k=0}^{\infty}\gamma^k
=
v_{\pi}(s)+\frac{c}{1-\gamma},
\]
for \(|\gamma|<1\).

Thus every state value is shifted by the same constant
\[
\boxed{v_c=\frac{c}{1-\gamma}.}
\]

Therefore,
\[
v'_{\pi}(s_1)-v'_{\pi}(s_2)
=
\left(v_{\pi}(s_1)+v_c\right)-\left(v_{\pi}(s_2)+v_c\right)
=
v_{\pi}(s_1)-v_{\pi}(s_2).
\]

So adding a constant to every reward does not change any relative ordering of states under
any policy. Only the intervals between rewards matter for comparing policies.

\section*{5}

Environment given:
\[
\begin{array}{c c c c c}
\toprule
\text{State} & \text{Action} & \text{Next State} & \text{Probability} & \text{Reward}\\
\midrule
0 & 0 & 0 & 1 & -1\\
0 & 1 & 1 & 1 & 1\\
1 & 0 & 1 & 1 & 1\\
1 & 1 & 0 & 1 & -1\\
\bottomrule
\end{array}
\]
with \(\gamma=0.1\).

The greedy action is \(a=1\) in state \(0\), and \(a=0\) in state \(1\).

\begin{enumerate}[label=\alph*)]
\item For two actions, an \(\varepsilon\)-greedy policy assigns
\[
\Pr(\text{greedy action})=1-\varepsilon+\frac{\varepsilon}{2}=1-\frac{\varepsilon}{2}.
\]
Given
\[
\pi(1|0)=\pi(0|1)=0.6,
\]
we have
\[
1-\frac{\varepsilon}{2}=0.6
\quad\Rightarrow\quad
\varepsilon=0.8.
\]

\item Under the given policy,
\[
V(0)=0.4[-1+0.1V(0)] + 0.6[1+0.1V(1)],
\]
\[
V(1)=0.6[1+0.1V(1)] + 0.4[-1+0.1V(0)].
\]

By symmetry, \(V(0)=V(1)=v\). Then
\[
v=0.4(-1+0.1v)+0.6(1+0.1v)=0.2+0.1v,
\]
so
\[
0.9v=0.2
\quad\Rightarrow\quad
v=\frac{2}{9}\approx 0.2222.
\]

Therefore,
\[
\boxed{V(0)=V(1)=\frac{2}{9}\approx 0.2222.}
\]

Now compute the action values:
\[
Q(0,0)=-1+0.1V(0)=-1+\frac{1}{45}=-\frac{44}{45}\approx -0.9778,
\]
\[
Q(0,1)=1+0.1V(1)=1+\frac{1}{45}=\frac{46}{45}\approx 1.0222,
\]
\[
Q(1,0)=1+0.1V(1)=\frac{46}{45}\approx 1.0222,
\]
\[
Q(1,1)=-1+0.1V(0)=-\frac{44}{45}\approx -0.9778.
\]

So
\[
\boxed{
Q(0,0)=-0.9778,\;
Q(0,1)=1.0222,\;
Q(1,0)=1.0222,\;
Q(1,1)=-0.9778.}
\]

\item The optimal Bellman equations are
\[
V^*(0)=\max\{-1+0.1V^*(0),\; 1+0.1V^*(1)\},
\]
\[
V^*(1)=\max\{1+0.1V^*(1),\; -1+0.1V^*(0)\}.
\]

Clearly the optimal actions are
\[
\boxed{\pi^*(0)=1,\qquad \pi^*(1)=0.}
\]

Under this policy,
\[
V^*(1)=1+0.1V^*(1)
\quad\Rightarrow\quad
V^*(1)=\frac{1}{0.9}=\frac{10}{9},
\]
and
\[
V^*(0)=1+0.1V^*(1)=1+\frac{1}{9}=\frac{10}{9}.
\]

Thus
\[
\boxed{V^*(0)=V^*(1)=\frac{10}{9}\approx 1.1111.}
\]

\item Using the given random numbers and the policy in part (a):

Trajectory starting from \(s_0=0\):
\[
0 \xrightarrow[a=0]{-1} 0 \xrightarrow[a=1]{1} 1 \xrightarrow[a=0]{1} 1.
\]
This uses \(0.94, 0.27, 0.30\).

Trajectory starting from \(s_0=1\):
\[
1 \xrightarrow[a=0]{1} 1 \xrightarrow[a=1]{-1} 0 \xrightarrow[a=1]{1} 1.
\]
This uses \(0.18, 0.75, 0.35\).

The returns are
\[
G_0^{(1)}=-1+0.1(1)+0.1^2(1)=-0.89,\qquad
G_1^{(1)}=1+0.1(1)=1.1,\qquad
G_2^{(1)}=1,
\]
\[
G_0^{(2)}=1+0.1(-1)+0.1^2(1)=0.91,\qquad
G_1^{(2)}=-1+0.1(1)=-0.9,\qquad
G_2^{(2)}=1.
\]

Using every-visit incremental Monte Carlo with \(\alpha=0.6\) and initial estimates \(0\):
\[
\boxed{V(0)\approx 0.77856,\qquad V(1)\approx -0.22560.}
\]

The action-value estimates are
\[
\boxed{
Q(0,0)=-0.534,\;
Q(0,1)=0.864,\;
Q(1,0)=0.786,\;
Q(1,1)=-0.540.}
\]

\item For off-policy TD control, use Q-learning:
\[
Q(s,a)\leftarrow Q(s,a)+\alpha\left[r+\gamma\max_{a'}Q(s',a')-Q(s,a)\right].
\]

Starting from all zeros, and treating the trajectory endpoint at \(t=3\) as terminal for the
last update, the final action values are
\[
\boxed{
Q(0,0)=-0.600,\;
Q(0,1)=0.840,\;
Q(1,0)=0.876,\;
Q(1,1)=-0.564.}
\]

\item For on-policy TD control, use SARSA:
\[
Q(s,a)\leftarrow Q(s,a)+\alpha\left[r+\gamma Q(s',a')-Q(s,a)\right].
\]

Start with \(Q(s,a)=0\) for all pairs and the initial policy
\[
\pi(1|0)=\pi(0|1)=0.5.
\]

Using the same random sequence from the start:
\[
0.94,\;0.27,\;0.30,\;0.18,\;0.75,
\]
the trajectory is
\[
0 \xrightarrow[a=0]{-1} 0 \xrightarrow[a=1]{1} 1 \xrightarrow[a=0]{1} 1
\xrightarrow[a=0]{1} 1 \xrightarrow[a=1]{-1} 0.
\]

Applying SARSA step-by-step gives
\[
\boxed{
Q(0,0)=-0.600,\;
Q(0,1)=0.600,\;
Q(1,0)=0.840,\;
Q(1,1)=-0.600.}
\]
\end{enumerate}

\section*{6}

The five 2D samples are
\[
(2,-2),\quad (0,4),\quad (-3,-1),\quad (3,-3),\quad (-2,2).
\]

Since the sample mean is
\[
\bm{\mu}=
\begin{bmatrix}
0\\
0
\end{bmatrix},
\]
the centered data matrix is the same as the original one:
\[
X=
\begin{bmatrix}
2 & 0 & -3 & 3 & -2\\
-2 & 4 & -1 & -3 & 2
\end{bmatrix}.
\]

The covariance matrix is
\[
\Sigma=\frac{1}{5}XX^{\top}
=
\begin{bmatrix}
5.2 & -2.8\\
-2.8 & 6.8
\end{bmatrix}.
\]

Its eigenvalues are
\[
\lambda_1\approx 8.9120,\qquad \lambda_2\approx 3.0880.
\]

Therefore the one-dimensional PCA direction is the eigenvector corresponding to \(\lambda_1\):
\[
\boxed{
\bm{u}_1=
\begin{bmatrix}
0.6022\\
-0.7983
\end{bmatrix}.}
\]

\begin{enumerate}[label=\alph*)]
\item One-dimensional PCA:
\[
\boxed{\bm{u}_1=\begin{bmatrix}0.6022\\-0.7983\end{bmatrix}.}
\]

\item The encoded 1D coordinates are
\[
z_i=\bm{u}_1^{\top}\bm{x}_i.
\]
Hence
\[
\boxed{
z=
\begin{bmatrix}
2.8011 & -3.1934 & -1.0082 & 4.2016 & -2.8011
\end{bmatrix}.}
\]

\item The reconstructed points are
\[
\hat{\bm{x}}_i=\bm{\mu}+\bm{u}_1 z_i.
\]
Therefore the reconstructed samples are approximately
\[
\boxed{
\begin{aligned}
\hat{\bm{x}}_1&=(1.6868,\,-2.2362),\\
\hat{\bm{x}}_2&=(-1.9230,\,2.5494),\\
\hat{\bm{x}}_3&=(-0.6072,\,0.8049),\\
\hat{\bm{x}}_4&=(2.5302,\,-3.3544),\\
\hat{\bm{x}}_5&=(-1.6868,\,2.2362).
\end{aligned}}
\]

\item The percentage of variance explained is
\[
\frac{\lambda_1}{\lambda_1+\lambda_2}
=
\frac{8.9120}{8.9120+3.0880}
\approx 0.7427.
\]
Therefore,
\[
\boxed{\text{explained variance ratio} \approx 74.27\%.}
\]
\end{enumerate}

\section*{7}

We have the data
\[
x=[1,2,3,9,10],
\]
and the initial parameters
\[
\pi_1=0.6,\qquad \pi_2=0.4,\qquad \mu_1=2.0,\qquad \mu_2=9.5,
\]
with \(\sigma_1^2=\sigma_2^2=1\).

For one EM iteration:

\textbf{E-step.} The responsibilities are
\[
\gamma_{i1}
=
\frac{\pi_1 \mathcal{N}(x_i\mid \mu_1,1)}
{\pi_1 \mathcal{N}(x_i\mid \mu_1,1)+\pi_2 \mathcal{N}(x_i\mid \mu_2,1)},
\qquad
\gamma_{i2}=1-\gamma_{i1}.
\]

Numerically,
\[
\begin{array}{c c c}
\toprule
x_i & \gamma_{i1} & \gamma_{i2}\\
\midrule
1  & 1.0000 & 0.0000\\
2  & 1.0000 & 0.0000\\
3  & 1.0000 & 0.0000\\
9  & 0.0000 & 1.0000\\
10 & 0.0000 & 1.0000\\
\bottomrule
\end{array}
\]
to four decimal places.

\[
N_1=\sum_i\gamma_{i1}\approx 3,\qquad N_2=\sum_i\gamma_{i2}\approx 2.
\]

\textbf{M-step.}
\[
\pi_1^{\text{new}}=\frac{N_1}{5}\approx 0.6,\qquad
\pi_2^{\text{new}}=\frac{N_2}{5}\approx 0.4.
\]

\[
\mu_1^{\text{new}}=\frac{\sum_i\gamma_{i1}x_i}{N_1}\approx \frac{1+2+3}{3}=2.0,
\]
\[
\mu_2^{\text{new}}=\frac{\sum_i\gamma_{i2}x_i}{N_2}\approx \frac{9+10}{2}=9.5.
\]

Therefore after one EM iteration,
\[
\boxed{
\pi_1^{\text{new}}\approx 0.6000,\quad
\pi_2^{\text{new}}\approx 0.4000,\quad
\mu_1^{\text{new}}\approx 2.0000,\quad
\mu_2^{\text{new}}\approx 9.5000.}
\]

This means that in practice the parameters remain essentially unchanged after one step because the two
clusters are already well separated.

\section*{8}

Interpret the listed joint probabilities in lexicographic order:
\[
P(X_0,X_1)=
\begin{bmatrix}
P(0,0)\\
P(0,1)\\
P(1,0)\\
P(1,1)
\end{bmatrix}
=
\begin{bmatrix}
\tfrac14\\
\tfrac12\\
\tfrac18\\
\tfrac18
\end{bmatrix}.
\]

Equivalently, the joint table is
\[
\begin{array}{c c c}
\toprule
 & X_1=0 & X_1=1\\
\midrule
X_0=0 & \tfrac14 & \tfrac12\\
X_0=1 & \tfrac18 & \tfrac18\\
\bottomrule
\end{array}
\]

\begin{enumerate}[label=\alph*)]
\item First compute the marginals:
\[
P(X_1=0)=\frac14+\frac18=\frac38,\qquad
P(X_1=1)=\frac12+\frac18=\frac58,
\]
\[
P(X_0=0)=\frac14+\frac12=\frac34,\qquad
P(X_0=1)=\frac18+\frac18=\frac14.
\]

Therefore
\[
P(X_0=0\mid X_1=0)=\frac{1/4}{3/8}=\frac23,\qquad
P(X_0=1\mid X_1=0)=\frac13,
\]
\[
P(X_0=0\mid X_1=1)=\frac{1/2}{5/8}=\frac45,\qquad
P(X_0=1\mid X_1=1)=\frac15.
\]

So
\[
\boxed{
P(X_0\mid X_1=0)=
\begin{bmatrix}
\tfrac23\\[2pt]
\tfrac13
\end{bmatrix},
\qquad
P(X_0\mid X_1=1)=
\begin{bmatrix}
\tfrac45\\[2pt]
\tfrac15
\end{bmatrix}.}
\]

Similarly,
\[
P(X_1=0\mid X_0=0)=\frac{1/4}{3/4}=\frac13,\qquad
P(X_1=1\mid X_0=0)=\frac23,
\]
\[
P(X_1=0\mid X_0=1)=\frac{1/8}{1/4}=\frac12,\qquad
P(X_1=1\mid X_0=1)=\frac12.
\]

Hence
\[
\boxed{
P(X_1\mid X_0=0)=
\begin{bmatrix}
\tfrac13\\[2pt]
\tfrac23
\end{bmatrix},
\qquad
P(X_1\mid X_0=1)=
\begin{bmatrix}
\tfrac12\\[2pt]
\tfrac12
\end{bmatrix}.}
\]

\item Use the random numbers in order.

Initialization from independent uniform distributions:
\[
r_1=0.94 \Rightarrow X_0^{(0)}=1,\qquad
r_2=0.27 \Rightarrow X_1^{(0)}=0.
\]

Now run Gibbs sampling for 3 iterations.

\textbf{Iteration 1}
\[
r_3=0.30,\quad X_1^{(0)}=0
\Rightarrow
P(X_0=0\mid X_1=0)=\frac23
\Rightarrow X_0^{(1)}=0.
\]
\[
r_4=0.18,\quad X_0^{(1)}=0
\Rightarrow
P(X_1=0\mid X_0=0)=\frac13
\Rightarrow X_1^{(1)}=0.
\]

\textbf{Iteration 2}
\[
r_5=0.75,\quad X_1^{(1)}=0
\Rightarrow X_0^{(2)}=1,
\]
because \(0.75>\tfrac23\).
\[
r_6=0.35,\quad X_0^{(2)}=1
\Rightarrow X_1^{(2)}=0,
\]
because \(0.35\le \tfrac12\).

\textbf{Iteration 3}
\[
r_7=0.61,\quad X_1^{(2)}=0
\Rightarrow X_0^{(3)}=0,
\]
because \(0.61\le \tfrac23\).
\[
r_8=0.42,\quad X_0^{(3)}=0
\Rightarrow X_1^{(3)}=1,
\]
because \(0.42>\tfrac13\).

So the Gibbs samples are
\[
\boxed{
(X_0^{(0)},X_1^{(0)})=(1,0),\quad
(X_0^{(1)},X_1^{(1)})=(0,0),\quad
(X_0^{(2)},X_1^{(2)})=(1,0),\quad
(X_0^{(3)},X_1^{(3)})=(0,1).}
\]
\end{enumerate}

\end{document}
